\appendices
\section{Reproducibility}
\label{app:repro}

\paragraph{Environment.}
The pinned versions are in \texttt{requirements.txt}: numpy 2.5.2, scipy 1.18.1,
scikit-learn 1.9.0, pandas 3.0.6, matplotlib 3.11.2, torch 2.14.0, transformers
5.16.1, sentence-transformers 6.0.1. The latency benchmark ran on Python 3.14.3;
every report also records the interpreter, platform and package versions it ran
with. The models are \texttt{gpt2} (perplexity),
\texttt{all-MiniLM-L6-v2} (embedding) and
\texttt{protectai/\allowbreak deberta-v3-base-\allowbreak prompt-injection-v2} (published classifier).

\paragraph{Seeds and splits.}
Seeds 42, 43 and 44 give the split hashes \texttt{\resHashA},
\texttt{\resHashB} and \texttt{\resHashC}; a rebuilt pool must reproduce them.
The held-out run writes a lock file \texttt{eval/results/final\_FINAL\_seed}$N$\texttt{.json}
per seed, and refuses to run again once it exists.

\paragraph{Commands.}
From the repository root:
{\footnotesize
\begin{verbatim}
pip install -r requirements.txt
python -m data.pool_datasets --beir-zip nq.zip
python -m eval.protocol --seed 42     # also 43, 44
python -m eval.final_test --seed 42 --dry-run
python -m eval.final_test --seed 42 --final
python -m eval.make_paper_tables --final
python paper/make_figures.py --final
\end{verbatim}}
The \texttt{--final} run is allowed once per seed, after the code is committed
and tagged. The dry run uses the validation role and writes
\texttt{final\_dryrun\_*} files.

\paragraph{What reproduces from the repository.}
The per-document scores, thresholds and reports of every reported run are
committed under \texttt{eval/results}, and the tables and figures of this paper are
generated from those files and from \texttt{latency\_benchmark.json} without
loading a model. Rebuilding the pooled document file needs downloads: the BIPIA
host documents, the PoisonedRAG targets and the BEIR Natural Questions corpus
(about 0.5\,GB), and the Hugging Face models above. The pooled file itself is not
in the repository.

\paragraph{Hardware.}
Latency (Table~\ref{tab:latency}) was measured on a Windows~11 laptop with an
Intel Core Ultra~5 125H (18 logical cores) and 16\,GB of RAM, plugged in, four
threads, 200 documents, three repeats, seed 42. Detector scores do not depend on
the hardware.

\paragraph{Licences and availability.}
BIPIA code and attack strings are MIT with component licences (WikiTableQuestions
and Stack Exchange CC BY-SA); the PoisonedRAG repository is MIT and Natural
Questions is CC BY-SA 3.0 (\texttt{docs/contamination.md}). Code, in-house
documents and scores are public at \url{https://github.com/tvxsha/pds}; the code
behind the reported runs is the tag \texttt{freeze-v1}.